%=======================================================================
%  FIGURE 1 - SYSTEM ARCHITECTURE  (spans both columns)
%  Edit the node text here to change the diagram. Everything is vector,
%  so it stays sharp at any zoom - never paste a screenshot instead.
%  Layout revised 2026-09-15: explicit coordinates so no two arrows share
%  a segment and no arrow passes through a box.
%=======================================================================
\begin{figure*}[t]
\centering
\resizebox{\textwidth}{!}{%
\begin{tikzpicture}[
  font=\sffamily\scriptsize,
  >={Stealth[length=2mm]},
  box/.style={rectangle, rounded corners=1.5pt, draw=black!70, line width=0.45pt,
              minimum height=7.5mm, align=center, inner xsep=2.2mm, fill=white},
  gen/.style={box, fill=blue!6},
  det/.style={box, fill=orange!12},
  act/.style={box, fill=green!10},
  dec/.style={diamond, aspect=2.1, draw=black!70, line width=0.45pt,
              align=center, inner sep=0.6mm, fill=yellow!14},
  term/.style={box, fill=black!8, rounded corners=3pt},
  lbl/.style={font=\sffamily\tiny, inner sep=0.8mm},
  ph/.style={font=\sffamily\bfseries\scriptsize}
]

% ---------------- generation + detection row ----------------
\node[term] (q)     at (0,0)    {Question $q$};
\node[gen]  (llm)   at (2.7,0)  {LLM decoding\\step $s_k$};
\node[gen]  (bnd)   at (5.5,0)  {Step boundary\\detector};
\node[det]  (hid)   at (8.4,0)  {Hidden state $\mathbf{h}_k$\\layer $\ell$};
\node[det]  (probe) at (11.1,0) {Linear probe\\$f_\theta$};
\node[det]  (risk)  at (14.0,0) {Risk $r_k$\\$=P(y{=}1\mid\mathcal{T}_{1:k})$};
\draw[->] (q)--(llm);
\draw[->] (llm)--(bnd);
\draw[->] (bnd)--(hid);
\draw[->] (hid)--(probe);
\draw[->] (probe)--(risk);

% answer is emitted by the decoder when generation ends
\node[term] (ans) at (2.7,1.55) {Final answer $a$};
\draw[->] (llm.north)--node[lbl,right]{answer emitted}(ans.south);

% offline training feeds the probe
\node[det] (train) at (11.1,1.55)
  {\textbf{Offline:} annotated traces $\to$ split \emph{by question}
   $\to$ fit $f_\theta$ $\to$ calibrate $\tau$ on corruption budget};
\draw[->, dashed, black!55] (train.south)--(probe.north);

% ---------------- decision ----------------
\node[dec] (thr)  at (14.0,-1.75) {$r_k \ge \tau$?};
\draw[->] (risk)--(thr);
\node[gen] (cont) at (5.5,-1.75) {Continue\\generation};
\draw[->] (thr)--node[lbl,above,pos=0.2]{no}(cont);
\draw[->] (cont.west) -| (llm.south);

% ---------------- selective mitigation ----------------
\node[act] (back) at (14.0,-3.35) {Backtrack to\\$b=\arg\max_{i\le k} r_i$};
\draw[->] (thr)--node[lbl,right]{yes}(back);
\node[act] (ver) at (11.4,-4.85) {Verify step $s_b$\\(cheap)};
\node[act] (reg) at (16.6,-4.85) {Regenerate from\\$s_{b-1}$ (costly)};
\draw[->] (back.west) -| node[lbl,above,pos=0.3]{$r_k<\tau_{\mathrm{hi}}$} (ver.north);
\draw[->] (back.east) -| node[lbl,above,pos=0.3]{$r_k\ge\tau_{\mathrm{hi}}$} (reg.north);
\node[dec] (imp) at (14.0,-6.3) {risk\\improved?};
\draw[->] (ver.south) |- (imp.west);
\draw[->] (reg.south) |- (imp.east);
\node[act] (rev)  at (12.2,-7.85) {Revert to $\mathcal{T}_{1:k}$};
\node[act] (keep) at (15.8,-7.85) {Keep revision $\mathcal{T}'$};
\draw[->] (imp.south west)--node[lbl,above left]{no}(rev.north);
\draw[->] (imp.south east)--node[lbl,above right]{yes}(keep.north);
\draw[->] (rev.west) -| (cont.south);
\draw[->] (keep.south) -- ++(0,-1.05) -| (cont.south);

% ---------------- phase brackets ----------------
\begin{scope}[on background layer]
  \node[fit=(hid)(probe)(risk)(thr), draw=orange!70, dashed, line width=0.6pt,
        rounded corners=2pt, inner sep=3.2mm] (p1) {};
  \node[fit=(back)(ver)(reg)(imp)(rev)(keep), draw=green!55!black, dashed, line width=0.6pt,
        rounded corners=2pt, inner sep=3.2mm] (p2) {};
\end{scope}
\node[ph, orange!70!black, anchor=south west] at ([yshift=0.8mm, xshift=1mm]p1.south west)
  {PHASE I --- DETECTION};
\node[ph, green!45!black, anchor=north west] at ([yshift=-0.8mm, xshift=1mm]p2.south west)
  {PHASE II --- SELECTIVE MITIGATION};

\end{tikzpicture}%
}
\caption{System architecture. \textbf{Phase I (detection)} evaluates a detector
$f_\theta$ at each step boundary to obtain a risk estimate $r_k$; the threshold
$\tau$ is calibrated offline against a corruption budget. \textbf{Phase II
(selective mitigation)} is entered only when $r_k \ge \tau$, escalating from
cheap step verification to costly regeneration, and reverting any revision that
does not reduce estimated risk. The default path is the \textsc{Continue} loop,
so the marginal cost on low-risk trajectories is one detector evaluation per step.}
\label{fig:arch}
\end{figure*}
